\documentclass[12pt]{article}
\title{Puncte Foarte Generale}
% Document
\begin{document}
\maketitle

Principalele puncte foarte generare ale proiectului sunt:
    \begin{itemize}
        \item Dezvoltarea unuia sau mai multor agenți, la alegere. Fie unul mare care sa le facă pe toate, fie mai mulți agenți micuți care să facă 1 singur task simplu(probabil alegerea potrivită este mai mulți agenți micuți care să facă 1 singur task.)
        \item Folosirea tehnicilor si/sau algoritmilor de Reinforcement Learning clasic și/sau Deep Reinforcement Learning, la alegere.
        \item Folosirea unui singur joc în care agenții vor desfășura acțiuni.
        \item Evaluarea agentilor prin metrici specifice RL/DRL.
        \item Comparații atât între agenți cât și cu literatura de specialitate. Eventual propuneri de îmbunătățire.
        \item Proiectul trebuie să fie rulabil și ușor de folosit.
        \item ( Optional ). Folosirea de teste pentru verificarea codului (Test Driven Development).
        \item Folosire predominant Python + pachete relevante.
        \item Daca este nevoie, dar într-o măsură limitată, folosirea și altor limbaje de programare. Dominant trebuie să rămână python + pachete.
        \item Dezvoltare interfață Python pentru controlul extern al jucătorilor.
    \end{itemize}
\end{document}